\chapter{Teeth whitening}
\index{Teeth whitening}

\medskip
\noindent\textit{Educational reference; always seek professional advice for individual care.}

\section*{Definition and scope}
Teeth whitening, also called tooth bleaching, is a cosmetic dental treatment that lightens the colour of the natural teeth and removes or reduces stains and discolouration. It is one of the most popular cosmetic dental procedures in many countries and around the world, chosen by people of all ages who want a brighter smile. Whitening works on the natural tooth structure, and it is important to understand that it does not change the colour of dental fillings, crowns, veneers, or dentures, which need to be matched separately if they are part of a smile makeover. There are two broad approaches: products applied at home, such as whitening strips, toothpastes, gels, and custom-fitted trays supplied by a dentist, and in-chair professional treatments performed in the dental surgery, often with a light or laser used to activate the bleaching agent. The active ingredient in almost all whitening products is a peroxide compound, most commonly hydrogen peroxide or carbamide peroxide, which breaks down stain molecules within the enamel. This chapter explains how whitening works, what results and side effects to expect, the difference between professional and at-home products, and how to choose a safe, effective approach. It also covers the situations in which whitening is not appropriate, such as teeth with untreated decay, gum disease, or allergy to peroxide, and when a dental opinion should be sought first. Throughout, the guiding principle is that a healthy mouth whitens best: whitening is a finish, not a fix, and it should never be used to mask an underlying dental problem.

\section*{Benefits and purpose}
The main benefit of whitening is cosmetic: most people find that lighter, brighter teeth improve their confidence and self-image, and many describe feeling more comfortable smiling and talking in social and professional settings. In a society where appearance matters, this modest but genuine improvement in self-esteem is the reason whitening remains one of the most requested procedures in general dental practice. Whitening can effectively address the common causes of tooth discolouration: surface stains from coffee, tea, red wine, and tobacco; age-related darkening of the dentine beneath the enamel; and the yellowing that occurs naturally as enamel thins over the years. For many people, whitening also reduces the need for more invasive cosmetic work, because a brighter natural tooth is less likely to prompt replacement with a crown or veneer. The purpose of professional supervision is to make the process safe: a dentist can confirm that the teeth and gums are healthy, choose the appropriate strength and duration of whitening, and manage sensitivity and other side effects. The purpose of this chapter is to give the reader the information needed to weigh up whitening options realistically, understand what it can and cannot achieve, and recognise the warning signs of an unsafe product or an unsupervised treatment.

\section*{How it works}
Tooth colour is determined by the enamel and the dentine beneath it, and stains sit either on the surface of the enamel or within the tooth structure itself. Surface stains are the brown and yellow films left by food, drink, and tobacco; they are removed by professional polishing and by the mild abrasives in whitening toothpastes. Internal discolouration is different: pigments become trapped within the enamel crystals and the dentine during the life of the tooth, sometimes strengthened by decades of staining, by some medicines, or by a developmental or traumatic event. It is this internal discolouration that true bleaching targets.

\begin{itemize}
    \item \textbf{Hydrogen peroxide:} the active bleaching agent in most professional and at-home products. It releases reactive oxygen molecules that penetrate the enamel and break down the pigmented chemical bonds responsible for discolouration, lightening the tooth from within
    \item \textbf{Carbamide peroxide:} a slower-acting compound that breaks down into hydrogen peroxide and urea when applied. It is common in at-home tray systems and custom-fitted whitening trays, because it releases its bleaching action gradually over several hours
    \item \textbf{Surface abrasives and polishing:} whitening toothpastes contain mild abrasives that remove surface stains only; they do not change the internal colour of the tooth and have limited effect on age-related yellowing
    \item \textbf{Light and laser activation:} some in-chair treatments use an activating light to accelerate the bleaching reaction. Evidence on the added benefit of the light is mixed, but it shortens chair time and many patients prefer the immediate result
    \item \textbf{Custom-fitted trays:} a dentist takes an impression and fabricates a tray that holds the gel against the teeth for a set time each day, usually for one to three weeks, giving controlled, even bleaching
    \item \textbf{Over-the-counter products:} strips, paint-on gels, and pre-filled trays apply a lower concentration of peroxide for shorter contact times, so results are slower, more variable, and depend heavily on correct application
\end{itemize}

The biology of the result explains a key limitation: whitening affects natural enamel only, and it is temporary. The teeth will gradually reabsorb stains from food, drink, and tobacco, so most people need a top-up treatment every six to twenty-four months to maintain the shade.

\section*{What to expect}
Before any whitening, a dentist should examine the teeth and gums, check for decay, cracked fillings, exposed tooth roots, and gum recession, and take note of the current shade. Teeth that are healthy whiten best, and treating any problems first avoids bleaching an already-damaged tooth. During an in-chair session, the gums and lips are protected with a barrier, the gel is applied, and the treatment runs for one to three sessions of roughly twenty to forty minutes each. The shade is checked as the procedure progresses. With home trays, the gel is worn for a period specified by the product, often thirty minutes to a few hours a day for one to three weeks. A realistic expectation is a lightening of several shades, usually described on a scale from a warm yellow to a cooler white. Results are usually noticeable after the first professional session and continue to improve over the following days as the teeth rehydrate and settle. Temporary tooth sensitivity to cold air, cold drinks, and occasionally sweets is the commonest side effect, affecting most people to some degree; it typically resolves within a few days to a week. The final shade cannot be predicted exactly, and very dark or severely stained teeth may not reach a brilliant white, a point that a reputable practitioner will raise honestly before treatment begins.

When the treatment is complete, the dentist will record the new shade and advise on maintenance. A common routine is a short top-up course, or the occasional use of a low-strength whitening toothpaste or strip, to hold the colour. Regular six-monthly cleans remove the surface stains that make teeth look duller than they are, and a professional clean alone can sometimes brighten the smile noticeably before any bleaching is even needed. It is also worth remembering that the perceived whiteness of teeth depends on context: a hygienic mouth with healthy gums and good alignment looks brighter than the raw shade reading alone would suggest.

\section*{When to seek professional advice}
Not every tooth that looks stained should be whitened, and several situations call for a dental assessment first. A professional consultation also allows the dentist to discuss realistic expectations, to photograph or shade-match the teeth as a baseline, and to recommend the treatment most likely to suit the individual's teeth, lifestyle, and budget. In many countries, high-strength whitening is legally restricted to dental professionals, so a dental visit is the safest starting point for anyone who wants a significant change in shade.

\begin{itemize}
    \item \textbf{Grey or brown single tooth:} a tooth that has died or been injured may need root canal treatment, and whitening will not restore it
    \item \textbf{Cavities or worn fillings:} bleach applied to decayed or leaking restorations can cause severe pain and aggravate the problem
    \item \textbf{Gum disease or recession:} exposed roots do not whiten well and are prone to sensitivity and irritation
    \item \textbf{Teeth with white patches, spots, or banding:} these marks from development, trauma, or excess fluoride often need specialist cosmetic techniques rather than simple bleaching
    \item \textbf{Pregnancy and breastfeeding:} whitening has not been studied in pregnancy and is generally deferred until after breastfeeding
    \item \textbf{Allergy:} anyone with a known reaction to peroxide should avoid whitening products
    \item \textbf{Teenagers:} younger teeth have larger pulp chambers and are more easily irritated; professional guidance is particularly important before age eighteen
\end{itemize}

People who grind their teeth, smokers, and heavy consumers of staining drinks should be advised that results will fade faster and maintenance will be needed more often.

\section*{Risks and cautions}
Whitening is safe when products are used correctly, but unsupervised or overused treatments carry genuine risks.

\begin{itemize}
    \item \textbf{Tooth sensitivity:} the commonest problem, usually temporary but occasionally lasting; it is worse with higher concentrations and longer contact times
    \item \textbf{Gum and soft-tissue burns:} peroxide left in contact with the gums causes white, painful patches; these settle, but they are uncomfortable and preventable with proper tray fit and barriers
    \item \textbf{Uneven or patchy results:} happens when trays fit poorly, the gel is applied thinly, or teeth have existing patches and restorations
    \item \textbf{Damage to tooth structure:} over-bleaching, or aggressive abrasive toothpastes used daily, can thin enamel and make teeth more translucent and more sensitive
    \item \textbf{Interactions with existing work:} crowns, veneers, and fillings do not change colour, so a whitened natural tooth can end up noticeably lighter than its restorations
    \item \textbf{Unregulated and fake products:} kits bought online, from market stalls, or through social media may contain unsafe concentrations of peroxide, illegal bleaching agents, or acids that damage enamel; local law restricts high-concentration products to dental professionals, and reputable practitioners will not sell unlabelled or imported gels
    \item \textbf{Oversight of underlying disease:} a stain can occasionally be a sign of decay, a dead nerve, or internal bleeding after injury, and whitening can delay the diagnosis of these problems
\end{itemize}

Anyone who experiences prolonged pain, swelling, or burning that does not settle should stop the treatment and see a dentist.

\section*{Tips and practical advice}
The following advice helps achieve the best result with the least harm.

\begin{itemize}
    \item See a dentist for an examination and a professional opinion before whitening, especially if there is any doubt about the health of the teeth and gums
    \item Choose a dentist-supervised treatment or a reputable local product, and follow the product instructions precisely; more is not better
    \item Use a toothpaste designed for sensitive teeth for a week or two before, during, and after whitening to reduce sensitivity
    \item Time whitening around events by starting a few weeks ahead, so sensitivity settles before the special occasion
    \item Maintain results with good brushing and flossing, regular professional cleans, and sensible limits on coffee, tea, red wine, and tobacco
    \item Avoid whitening during pregnancy and breastfeeding, and avoid treating young children's teeth without professional advice
    \item Be wary of very cheap kits, "guaranteed" brilliant-white claims, and products sold outside pharmacies and dental practices
    \item Revisit the dentist if the teeth become persistently sensitive, patchy, or if pain develops
\end{itemize}

\section*{Sources and further reading}
The following organisations provide reliable information about teeth whitening and dental health.

\begin{itemize}
    \item Australian Dental Association. \textit{Teeth whitening: what you need to know.} https://www.ada.org.au/
    \item healthdirect Australia. \textit{Teeth whitening and teeth whitening products.} https://www.healthdirect.gov.au/
    \item NHS. \textit{Teeth whitening.} https://www.nhs.uk/
    \item Mayo Clinic. \textit{Teeth whitening: what are your options?} https://www.mayoclinic.org/
    \item Australian Department of Health and Aged Care. \textit{Oral health information for the public.} https://www.health.gov.au/
    \item Therapeutic Goods Administration. \textit{Regulation of teeth whitening products.} https://www.tga.gov.au/
\end{itemize}

\clearpage
